% 作业 1：GPU 与 GPU 编程
% 编译：xelatex -interaction=nonstopmode assignment01.tex（跑两遍）
\documentclass[11pt,a4paper,fontset=none]{ctexart}

\setCJKmainfont[BoldFont={Noto Serif CJK SC Bold}]{Noto Serif CJK SC}
\setCJKsansfont[BoldFont={Noto Sans CJK SC Bold}]{Noto Sans CJK SC}
\setCJKmonofont{Noto Sans CJK SC}
\usepackage{geometry}
\geometry{margin=2.3cm}

\usepackage{amsmath,amssymb}
\usepackage{booktabs,array,multirow}
\usepackage{enumitem}
\setlist{nosep,leftmargin=2em}
\usepackage{xcolor}
\definecolor{accent}{RGB}{28,60,120}
\definecolor{codebg}{RGB}{248,248,246}
\usepackage[most]{tcolorbox}
\usepackage{listings}
\definecolor{strgreen}{RGB}{21,101,42}
\lstset{
  basicstyle=\ttfamily\small,
  backgroundcolor=\color{codebg},
  frame=single, rulecolor=\color{gray!40}, framerule=0.4pt,
  breaklines=true, columns=fullflexible, keepspaces=true,
  showstringspaces=false,
  language=bash,
  morekeywords={make,nvcc,uv,pytest,python,cd},
  commentstyle=\color{gray!130}\itshape,
  keywordstyle=\color{accent}\bfseries,
  stringstyle=\color{strgreen},
  xleftmargin=4pt, xrightmargin=4pt,
}
\usepackage{hyperref}
\hypersetup{colorlinks, linkcolor=accent, urlcolor=accent}
\emergencystretch=3em

% 参考资料框
\newtcolorbox{reading}{enhanced, breakable,
  colback=gray!5, colframe=gray!50, boxrule=0.5pt, arc=1.5pt,
  left=7pt, right=7pt, top=5pt, bottom=5pt,
  title={参考资料}, fonttitle=\bfseries\sffamily\small,
  coltitle=black, attach title to upper={\quad}}

% 压轴题框
\newtcolorbox{capstone}[1]{enhanced, breakable,
  colback=accent!4, colframe=accent!70, boxrule=0.7pt, arc=1.5pt,
  left=7pt, right=7pt, top=5pt, bottom=5pt,
  title={#1}, fonttitle=\bfseries\sffamily}

% 回望框
\newtcolorbox{lookback}{enhanced, breakable,
  colback=gray!5, colframe=gray!50, boxrule=0.5pt, arc=1.5pt,
  left=7pt, right=7pt, top=5pt, bottom=5pt,
  title={Editor's Note}, fonttitle=\bfseries\sffamily\small,
  coltitle=black, attach title to upper={\quad},
  before upper app={\setlength{\parskip}{0.5em}}}

% 题目标题：\prob{编号}{TYPE}；选做题写 \prob[Optional]{编号}{TYPE}，必做不标注；
% 题目主要涉及的文件用末尾的可选参数标在标题行右侧：\prob{编号}{TYPE}[相对路径]
\NewDocumentCommand{\prob}{O{}mmo}{\par\medskip\noindent
  {\bfseries prob #2}\;{\small\sffamily（#3\IfBlankF{#1}{\,·\,#1}）}%
  \IfValueT{#4}{\hfill{\footnotesize\ttfamily\color{black!60}#4}}\par\nobreak\smallskip\noindent\ignorespaces}

\newcommand{\guide}{CUDA Programming Guide}

\begin{document}

\begin{center}
  {\LARGE\bfseries 作业 1：GPU 与 GPU 编程}\\[6pt]
  {\sffamily Weiming HPC Training Camp $\times$ LCPU AI Infra Seminars \; · \; Session 1}
\end{center}

\vspace{0.5em}
\hrule
\vspace{1em}

\section*{Preface}

此次作业为 Session 1 的配套练习，内容覆盖 \guide{}\footnote{\url{https://docs.nvidia.com/cuda/cuda-programming-guide/}，2025 年重排的新版。下文引用记作 Guide。} 第一、第二部分的全部章节。

共有七种题型：CONCEPT 概念题、FILL-IN 填空题、MODIFY 改造题、DEBUG 修 bug 题、EXPERIMENT 实验题、HANDS-ON 动手题、FROM-SCRATCH “从零实现”题。标 Optional 的为选做题。涉及代码编辑的题目会在标题行右侧标出相应代码文件的路径（相对 \texttt{assignment01/}）。FROM-SCRATCH 题会配套判测脚本。

代码在仓库 \texttt{assignment01/} 目录，环境配置见其中的 \texttt{README.md}。

AI Policy: \texttt{CLAUDE.md} 或 \texttt{AGENTS.md}。核心原则是“AI 可以帮你理解，但不能替你实现”。

\subsection*{Content}

\begin{center}\small
\begin{tabular}{@{}llll@{}}
\toprule
Module & 主题 & 对应阅读 & 代码位置 \\
\midrule
0 & 环境准备 & \texttt{README.md} & \texttt{cuda/m0\_env/} \\
1 & 为什么要用 GPU & Guide 1.1、1.2.1--1.2.2 & \texttt{cuda/m1\_why\_gpu/} \\
2 & 第一个 CUDA 程序 & Guide 2.1 全章、2.6 & \texttt{cuda/m2\_first\_kernel/} \\
3 & SIMT 执行 & Guide 2.3.1--2.3.2、1.2.2.2 & \texttt{cuda/m3\_simt/} \\
4 & 存储空间 & Guide 2.3.3--2.3.5、2.3.7、1.2.3 & \texttt{cuda/m4\_memory/} \\
5 & 计时与异步初步 & Guide 2.5 引言、1.2.3.3 & \texttt{cuda/m5\_async/} \\
6 & Tile 视角 & Guide 1.2.2.3、2.4、2.2 & （阅读为主） \\
7 & TileLang 与 Triton & TileLang 文档、Triton 教程 & \texttt{kernels/} + \texttt{tests/} \\
8 & 平台与编译 & Guide 1.3、2.1.1、2.7 & （复用 \texttt{cuda/m0\_env}） \\
Bonus & matmul & — & \texttt{cuda/bonus/}、\texttt{kernels/} \\
\bottomrule
\end{tabular}
\end{center}

\setcounter{section}{-1}

% ================================================================
\section{环境准备}

先把环境跑通——编译一个 minimal 的 CUDA 程序，确认工具链和卡都能用；再把这块卡的参数查一遍，后面的 Module 会反复用到。环境配置见 \texttt{README.md} 。

\begin{reading}
\texttt{assignment01/README.md}；查硬件参数用 Guide 第五部分附录 Compute Capabilities。
\end{reading}

\prob{0.1}{HANDS-ON}[cuda/m0\_env/01\_hello.cu]
编译并运行第一个程序，它启动 4 个 block、每个 block 8 个线程。

\begin{lstlisting}
cd assignment01/cuda
make run/m0_env/01_hello
\end{lstlisting}

看看输出的结果，建议至少运行 5 次，可以留意输出各行顺序的变化。源码可以先不细读，Module 8 会回头看 \texttt{nvcc} 对它做了什么。

\prob{0.2}{FILL-IN}[cuda/m0\_env/02\_device\_query.cu]
\texttt{02\_device\_query.cu} 的五个空都是 \texttt{cudaDeviceProp} 的字段名，对照 CUDA Runtime API 文档补全，然后编译运行。

\begin{lstlisting}
cd assignment01/cuda
make run/m0_env/02_device_query
\end{lstlisting}

把你这块卡的参数填进下表，并对照 Guide 附录 Compute Capabilities 里对应架构的表，核对 warp 大小和每 SM 最大线程数。

\begin{center}\small
\begin{tabular}{@{}ll@{}}
\toprule
项目 & 你的卡 \\
\midrule
型号 / compute capability & \\
SM 数量 & \\
warp 大小 & \\
shared memory / block & \\
最大常驻线程 / SM & \\
显存总量 & \\
\bottomrule
\end{tabular}
\end{center}

\begin{lookback}
本模块两道题看上去只是热身，但 prob 0.2 的那张表后面会反复被引用：Module 3 讨论 warp 要用 warp 大小，Module 4 手算驻留 block 数要用“shared memory / SM”与“最大常驻线程 / SM”。Keep it at hand.

另外，我们将会逐渐领悟到：一段在 GPU 上运行的代码，其性能与硬件型号是强绑定的。同一段代码在不同 compute capability 上的性能可以相差数倍，某些优化只在某几代 compute capability 上奏效。此后每记录一个耗时，请将硬件型号与 compute capability 一同记下。
\end{lookback}

% ================================================================
\section{为什么要用 GPU}

GPU 与 CPU 的根本区别在设计目标——CPU 为单线程延迟服务，GPU 为总吞吐服务。此 Module 关注延迟/吞吐、SIMD/SIMT 等相关概念。

\begin{reading}
Guide 1.1、1.2.1--1.2.2；LCPU 讲义\footnote{\url{https://github.com/interestingLSY/CUDA-From-Correctness-To-Performance-Code/blob/master/lecture.md}} Part 0。
\end{reading}

\prob{1.1}{CONCEPT}
判断对错，可以顺带补一句理由。
\begin{enumerate}[label=(\alph*)]
  \item 一块标称 100 TFLOPS 的 GPU，执行单条指令的延迟一定低于 5 GHz 的 CPU。
  \item HBM 的“高带宽”指大块连续访问时的吞吐，零散的随机访问达不到标称值。
  \item 严格串行的迭代算法（每步依赖上一步的结果），即使换一块算力更强的 GPU 也快不了多少。
  \item “算力 1000 TFLOPS”意味着每次运算的延迟是 $10^{-15}$ 秒。
\end{enumerate}

\prob{1.2}{CONCEPT}
Session 1 讲座里提过“N 方过百万”这个例子。总计算量 $10^{12}$ FLOP 在当代 GPU 上的运算时间大概是毫秒级，那为什么一个严格在线的串行算法仍然做不到几秒内跑完？（从“延迟”和“吞吐”的角度考虑）

\prob{1.3}{CONCEPT}
补全下表（thread 一行已填好作为示例）

\begin{center}\small
\begin{tabular}{@{}p{5.2em}p{9em}p{7.5em}p{7.5em}p{8em}@{}}
\toprule
执行层次 & 软件含义 & 对应硬件 & 直接可用的存储 & 同步与通信手段 \\
\midrule
thread & kernel 的最小执行单位 & 计算单元上的一个 lane & 自己的寄存器 & （自身天然有序） \\
warp &  &  &  &  \\
block / CTA &  &  &  &  \\
grid &  &  &  &  \\
\bottomrule
\end{tabular}
\end{center}

\prob{1.4}{CONCEPT}
SIMD 与 SIMT 的区别？另：判断正误——Nvidia GPU 在 Volta 之后每个线程有独立的 program counter，所以 branch divergence 不再有性能代价。

\prob{1.5}{EXPERIMENT}[cuda/m1\_why\_gpu/01\_scaling.cu]
同一个向量加法，四种配置各跑一遍。

\begin{lstlisting}
cd assignment01/cuda
make run/m1_why_gpu/01_scaling
\end{lstlisting}

将数据填入下表，并回答：(a)~GPU 单线程为什么比 CPU 慢这么多？(b)~从单 block 到铺满 grid 的提速，说明 GPU 加速计算靠的是什么？

\begin{center}\small
\begin{tabular}{@{}lll@{}}
\toprule
配置 & 耗时 (ms) & ns / 元素 \\
\midrule
CPU 单线程 & & \\
GPU \texttt{<<<1,1>>>} & & \\
GPU \texttt{<<<1,256>>>} & & \\
GPU 铺满 grid & & \\
\bottomrule
\end{tabular}
\end{center}

\prob[Optional]{1.6}{FROM-SCRATCH}[kernels/simt\_sim.py]
SIMT Simulator —— 写一个 warp 的执行模拟器：32 个 lane、共享控制流、带 mask 的分支执行与汇合。请在 \texttt{kernels/simt\_sim.py} 内完成（docstring 里面有具体实现要求）。工作量大概是 50 行以内的 Python。不需要 GPU。注：本题为 Module 3 的实验结果提供理论对照。判测命令如下：

\begin{lstlisting}
cd assignment01
uv run pytest tests/test_simt_sim.py
\end{lstlisting}

\begin{lookback}
延迟与吞吐的分野根植于设计取舍，而非工艺差异。把“同一笔晶体管预算”投向不同的方向：一边是乱序执行、分支预测与庞大的 cache，另一边则是更多计算单元和更宽的访存带宽。prob 1.5 用四档配置让这个判断变得可度量，prob 1.1 与 prob 1.4 分别从两侧划出了它的边界，选做的 prob 1.6 则预先准备了一份理论标尺，留待 Module 3 的实测去校准。

更值得记住的是 prob 1.2 得出的那条结论：无论可用吞吐有多高，它始终绕不过一条串行依赖链。判断某个任务是否值得搬上 GPU，最先看的不是总计算量，而是依赖“图”的宽度——Amdahl 定律在 GPU 场景下的变体大抵如此。至于 prob 1.5 的那份向量加法，它几乎榨不出这块卡真正的算力，真正的瓶颈藏在另一个地方；Module 4 会为你把它量化出来。
\end{lookback}

% ================================================================
\section{第一个 CUDA 程序}

此 Module 把一个 CUDA 程序的完整生命周期过一遍，覆盖 Guide 2.1 的每个小节。

\begin{reading}
Guide 2.1 全章（主要出处）、2.6（浏览 unified memory 相关小节）；LCPU 讲义 Part 1。
\end{reading}

\prob{2.1}{FILL-IN}[cuda/m2\_first\_kernel/01\_vector\_add.cu]
请填空：\texttt{m2\_first\_kernel/01\_vector\_add.cu} ，具体要求见代码注释。判测命令如下：

\begin{lstlisting}
cd assignment01/cuda
make run/m2_first_kernel/01_vector_add
\end{lstlisting}

\prob{2.2}{CONCEPT}
为下列五个场景选择正确的修饰符（如 \texttt{\_\_global\_\_} 等）。
\begin{enumerate}[label=(\alph*)]
  \item 在 GPU 上执行、由 CPU 侧启动的 kernel 函数。
  \item 只会被 kernel 调用的辅助函数。
  \item host 和 device 代码都要调用的小工具函数。
  \item 整个 kernel 运行期间不变、所有线程都要读的系数表。
  \item block 内线程共享的暂存数组。
\end{enumerate}

\prob{2.3}{MODIFY}[cuda/m2\_first\_kernel/02\_vector\_add\_um.cu]
\texttt{02\_vector\_add\_um.cu} 代码完整，但目前是“显式内存管理”的版本。改之前先按原样编译运行一次，记下耗时——这一版会被你的改动覆盖掉，下面 (b) 要拿它做对照。然后按文件头的说明改成 unified memory 版（\texttt{cudaMallocManaged}），并保持文件头写明的计时窗口不变。

\begin{lstlisting}
cd assignment01/cuda
make run/m2_first_kernel/02_vector_add_um
\end{lstlisting}

然后请回答如下问题：(a)~kernel 启动之后、CPU 读结果之前，为什么必须有一次同步？在原先的版本里这次同步发生在哪个调用里？(b)~对比两版“搬运 + kernel + 读回”的耗时，分析差距的原因（谁快谁慢都有可能，与使用的卡有关）。

\prob{2.4}{CONCEPT}
判断对错，可以顺带补一句理由。
\begin{enumerate}[label=(\alph*)]
  \item \texttt{vectorAdd<<<...>>>(...)} 这条语句返回时，kernel 一定已经执行完毕。
  \item 同一个 stream 里，\texttt{cudaMemcpy}（device 到 host）会等它前面的 kernel 全部完成后才开始拷贝。
  \item kernel 内部的非法访存，会在启动语句处同步地报出来。
\end{enumerate}

\prob{2.5}{DEBUG}[cuda/m2\_first\_kernel/03\_bug\_launch.cu]
修 bug：\texttt{03\_bug\_launch.cu}，详细内容见相关文件。

\begin{lstlisting}
cd assignment01/cuda
make run/m2_first_kernel/03_bug_launch
\end{lstlisting}

\prob{2.6}{FILL-IN}[cuda/m2\_first\_kernel/04\_matrix\_add.cu]
\texttt{04\_matrix\_add.cu} 用二维的 block 和 grid 处理 $1000 \times 700$ 的矩阵，请填空实现矩阵加法。测试命令：

\begin{lstlisting}
cd assignment01/cuda
make run/m2_first_kernel/04_matrix_add
\end{lstlisting}

\prob{2.7}{MODIFY}[cuda/m2\_first\_kernel/05\_grid\_stride.cu]
\texttt{05\_grid\_stride.cu} 的 launch 被固定成 \texttt{<<<64, 256>>>}，线程总数远小于 $n$，当前 FAIL。在 launch 配置不变的前提下，把 kernel 改成 grid-stride loop，让任意 $n$ 都能 PASS。

\begin{lstlisting}
cd assignment01/cuda
make run/m2_first_kernel/05_grid_stride
\end{lstlisting}

然后请回答——这种写法的价值在哪里？launch 只有 16384 个线程时，性能上要付出什么代价？

\prob{2.8}{EXPERIMENT}[cuda/m2\_first\_kernel/06\_whoami.cu]
运行下面的程序两三次，观察 16 个 block 打印输出的先后。

\begin{lstlisting}
cd assignment01/cuda
make run/m2_first_kernel/06_whoami
\end{lstlisting}

(a)~顺序由谁决定？(b)~程序的正确性可以依赖 block 的执行顺序吗？这条限制和 Guide 1.1 说的 scalable programming model 有什么关系？

\begin{capstone}{prob 2.9（FROM-SCRATCH）：SAXPY \hfill {\footnotesize\ttfamily cuda/m2\_first\_kernel/saxpy.cu}}
在 \texttt{m2\_first\_kernel/} 下写出完整的 CUDA 程序 \texttt{saxpy.cu}，实现 $y \leftarrow 2.0 \cdot x + y$（单精度）。要求如下:
\begin{itemize}
  \item 不允许 include \texttt{common.h}。错误检查宏和 \texttt{cudaEvent} 计时都要自己写一遍。
  \item 命令行用法：\texttt{./saxpy <n>}，$n$ 是元素个数。输入数据按固定公式生成（都是 float）：\texttt{x[i] = ((i \% 2048) - 1024) * 0.5f}，\texttt{y[i] = (i \% 1024) - 512}。
  \item kernel 算完把 $y$ 拷回 host，用 double 累加所有 \texttt{y[i]}，输出一行 \texttt{SUM=<总和>}（用 \texttt{printf("SUM=\%.0f\textbackslash n", s)} 这样的格式，同一行里可以再带上 $n$ 和 kernel 毫秒数），SUM 结果将用于对拍检验程序正确性，exit code 应为 0。
  \item $n = 0$ 时输出 \texttt{SUM=0}，exit code 为 0（0 个 block 的 kernel launch 是非法的，特判即可）。
\end{itemize}
\end{capstone}

判测脚本覆盖 $n \in \{0, 1, 31, 1024, 1025, 2^{20}, 2^{20}{+}3\}$，命令如下。

\begin{lstlisting}
cd assignment01/cuda/m2_first_kernel
./judge_saxpy.sh saxpy.cu
\end{lstlisting}

注：SAXPY 即 Single-precision A$\cdot$X Plus Y。

\begin{lookback}
CUDA C++ 在 C++ 之上添加的语法扩展其实相当克制：几种函数与变量修饰符、一个 kernel 启动语法，以及一组内存管理 API。真正需要你花时间去适应的是另外两件事——host 与 device 各自独立推进的时间线（prob 2.4），以及数据此刻落在谁的地址空间里，你得“显式”地负责（prob 2.3）。作为本模块收尾的 prob 2.9 把这些概念连同索引、边界检查与错误处理一起收进一个不到百行的程序，prob 5.1 还会回过头来检验其中的计时部分。

prob 2.7 和 prob 2.8 值得多花一些时间。你为这两道题写下的理由——一条关于启动配置与问题规模的关系，另一条关于 block 之间的执行顺序——在后面的内容里会以不同面貌反复出现。Module 7 介绍的 Triton 与 TileLang 仍然建立在同一条 block 无序执行的假设之上，只不过不再需要你亲手把它写出来。
\end{lookback}

% ================================================================
\section{SIMT 执行}

SIMT 给每个线程“独立执行”的表象，硬件却按 32 线程一组共享 instruction fetch 和 issue（取指和发射）。此模块三个实验的主题分别为“divergence 的代价”、“\texttt{\_\_syncthreads} 的必要性”、“block 之间为什么无法同步”。

\begin{reading}
Guide 2.3.1--2.3.2、1.2.2.2。3.3、3.5 要用 shared memory，用法见 Guide 2.3.3.2 与 1.2.3.2（Module 4 会系统地读这一块）；cooperative groups 见 Guide 2.3.6（只需浏览）。
\end{reading}

\prob{3.1}{CONCEPT}
设 \texttt{blockDim = (8, 8, 1)}。
\begin{enumerate}[label=(\alph*)]
  \item \texttt{threadIdx = (3, 5, 0)} 的线性编号是多少？它在第几个 warp、warp 内第几个 lane？
  \item 这个 block 一共占多少个 warp？
  \item 若 \texttt{blockDim = (33, 1, 1)}，占几个 warp？这样配置浪费在哪里？
\end{enumerate}

\prob{3.2}{EXPERIMENT}[cuda/m3\_simt/01\_divergence.cu]
\texttt{m3\_simt/01\_divergence.cu} 的两个 kernel 每线程计算量相同，分支划分不同——一个按 thread 编号的奇偶分（同一个 warp 里一半一半），一个按 warp 边界对齐分。请先预测一下哪个版本运行会更快一点，大概快几倍，然后运行验证：

\begin{lstlisting}
cd assignment01/cuda
make run/m3_simt/01_divergence
\end{lstlisting}

请解释实测比值，并回答——若两个分支的计算量一大一小，按 thread 编号奇偶分的 kernel 和按 warp 边界对齐分的 kernel 的运行时间分别由什么决定？

\prob{3.3}{EXPERIMENT}[cuda/m3\_simt/02\_sync\_matters.cu]
\texttt{02\_sync\_matters.cu} 让每个 block 用 shared memory 把自己的 256 个元素倒序。请按文件开头的注释内容进行实验。

\begin{lstlisting}
cd assignment01/cuda
make run/m3_simt/02_sync_matters
\end{lstlisting}

(a)~为什么注释掉 sync 后代码不能正确地运行？(b)~(Optional) 注释掉 sync 后，翻转后的数组错的位置比较随机，但是有些位置一直是对的，试解释原因。（tip: 算一算 $t$ 与 $255-t$ 有没有可能落在同一个 warp）

\prob{3.4}{CONCEPT}
\texttt{\_\_syncthreads} 只能同步本 block 内的 threads，那需要全 grid 同步时，标准做法是什么？

\begin{capstone}{prob 3.5（FROM-SCRATCH）：block 内归约 \hfill {\footnotesize\ttfamily cuda/m3\_simt/03\_reduce.cu}}
在 \texttt{03\_reduce.cu} 里从零实现两个求和归约 kernel（判测与计时的代码已经写好）。两个 kernel 的 contract 见文件头。PASS 后，试解释实测性能差距的原因。

（Optional）基于两点事实——(a) \texttt{\_\_shfl\_down\_sync} 是 warp 内寄存器级别的线程间数据交换指令，自带同步效果且延迟极小；(b) 归约到最后 32 个元素后，活跃线程若都落在同一个 warp 里，就不再需要 \texttt{\_\_syncthreads}（可以想想为什么）——据此试写出第三版优化后的 kernel。测试时只会跑前两版，第三版自己在 \texttt{main} 里照着加一次 \texttt{run\_one} 调用即可。
\end{capstone}

\begin{lstlisting}
cd assignment01/cuda
make run/m3_simt/03_reduce
\end{lstlisting}

\begin{lookback}
SIMT 让你以单线程的风格编写并行代码，但它只承诺正确性语义，不负责性能语义。本模块的题目正好分别对应这层抽象漏出来的三个口子：warp 才是真实的调度单位（prob 3.2），block 内的并行需要显式同步才能看到彼此的写入（prob 3.3），而 block 之间压根没有提供同步方法（prob 3.4）。

``For performance reasoning, you always need to fall back to the warp level.''

prob 3.5 抛出的结论更让人不适：两个版本的加法次数完全相同，耗时却能相差超过一倍。算法复杂度在这里不再是性能的充分描述——这大概是 GPU 编程与经典算法课之间最根本的分歧。选做的第三版则指向另一条路径：不再把 warp 当作需要小心绕开的硬件细节，而是直接把它当作可调用的编程接口。warp-level primitive 与 cooperative groups（Guide 2.3.6）正是这个方向的官方解答。
\end{lookback}

% ================================================================
\section{存储空间}

数据的存储位置和迁移速度极大地影响着 kernel 的速度。此 Module 旨在对“存储如何影响性能”建立起基础认知。

\begin{reading}
Guide 2.3.3--2.3.5、2.3.7、1.2.3。
\end{reading}

\prob{4.1}{CONCEPT}
补全下表：

\begin{center}\small
\begin{tabular}{@{}lp{8em}p{7em}p{6em}p{8em}@{}}
\toprule
空间 & 谁可见 & 生命周期 & 片上 / 片外 & 谁管理 \\
\midrule
register & 单个线程 & 线程 & 片上 & 编译器 \\
local &  &  &  &  \\
shared &  &  &  &  \\
global &  &  &  &  \\
constant &  &  &  &  \\
L1 / L2 cache &  &  &  &  \\
\bottomrule
\end{tabular}
\end{center}

\prob{4.2}{FILL-IN}[cuda/m4\_memory/01\_stencil.cu]
请填空：\texttt{m4\_memory/01\_stencil.cu} ，具体要求见代码注释。测试命令：

\begin{lstlisting}
cd assignment01/cuda
make run/m4_memory/01_stencil
\end{lstlisting}

\prob{4.3}{MODIFY}[cuda/m4\_memory/02\_constant\_coeff.cu]
\texttt{02\_constant\_coeff.cu} ：按文件头的说明进行修改。修改完后测试：

\begin{lstlisting}
cd assignment01/cuda
make run/m4_memory/02_constant_coeff
\end{lstlisting}

结果会包含两个版本的耗时（差距可能极小，甚至测不出来性能收益——想想为什么），回答 constant cache 真正的优势在哪种访问模式。

\prob{4.4}{CONCEPT}
判断对错，可以顺带补一句理由。
\begin{enumerate}[label=(\alph*)]
  \item local memory 的“local”指作用域私有，它实际上在片外显存里。
  \item 对数组用运行期才知道的下标做索引，可能迫使它被放进 local memory。
\end{enumerate}

\prob{4.5}{FILL-IN}[cuda/m4\_memory/03\_histogram.cu]
请填空：\texttt{03\_histogram.cu} ，具体要求见代码注释。测试命令：

\begin{lstlisting}
cd assignment01/cuda
make run/m4_memory/03_histogram
\end{lstlisting}

\prob{4.6}{MODIFY}[cuda/m4\_memory/04\_histogram\_priv.cu]
\texttt{04\_histogram\_priv.cu} ：请按要求修改代码。改完测试指令：

\begin{lstlisting}
cd assignment01/cuda
make run/m4_memory/04_histogram_priv
\end{lstlisting}

测试结果包含与 naive 实现相比较的耗时、吞吐。试解释提速来自哪里。

\prob{4.7}{EXPERIMENT}[cuda/m4\_memory/05\_bandwidth.cu]
运行实验，并根据实验数据填表。

\begin{lstlisting}
cd assignment01/cuda
make run/m4_memory/05_bandwidth
\end{lstlisting}

观察数据变化趋势，并简析趋势的成因。

\begin{center}\small
\begin{tabular}{@{}lcccccc@{}}
\toprule
stride & 1 & 2 & 4 & 8 & 16 & 32 \\
\midrule
GB/s & & & & & & \\
\bottomrule
\end{tabular}
\end{center}

\prob{4.8}{EXPERIMENT}[cuda/m4\_memory/06\_occupancy.cu]
\texttt{06\_occupancy.cu}，详细内容见相关文件。实验命令如下：

\begin{lstlisting}
cd assignment01/cuda
make run/m4_memory/06_occupancy
\end{lstlisting}

记录实验数据：

\begin{center}\small
\begin{tabular}{@{}lcccccc@{}}
\toprule
shared memory / block (KB) & & & & & & \\
理论驻留 block / SM & & & & & & \\
occupancy & & & & & & \\
实测带宽 (GB/s) & & & & & & \\
\bottomrule
\end{tabular}
\end{center}

请回答：(a)~用程序开头打印的“shared memory / SM”和“最大常驻线程 / SM”，手算其中一个的驻留 block 数和 occupancy，和 API 的结果对照。(b)~带宽为什么随 occupancy 下降？用“延迟隐藏需要足够多的常驻 warp”组织你的解释。(c)~表中带宽随 occupancy 单调下降，但明显不成正比——从 100\% 到 75\% 带宽掉了多少？从 37.5\% 到 12.5\% 又掉了多少？试解释这个差别。

\begin{lookback}
本模块的八道题围绕同一个问题展开：数据放在哪一层，以什么模式去取。prob 4.1 给出一幅静态的存储层次地图，prob 4.2 与 4.3 在同一段 stencil 上切换存储位置，prob 4.5 与 4.6 在同一个直方图统计上做同样的变换，prob 4.7 和 4.8 则分别量化访问模式与可并发的请求数量。

有三点值得单独记下。其一，优化只在真正的瓶颈处生效。prob 4.3 是一个刻意安排的负结果：教科书罗列的手段用在了不构成瓶颈的地方，收益为零。避免这类空转，需要先估算再动手，估算可以依赖 arithmetic intensity 与 roofline 模型，后面的 session 会展开。其二，你在 prob 4.8 (b) 写下的那句解释有一个正式的名称——Little 定律：维持某一吞吐量所需的在途请求数，等于吞吐与延迟的乘积。其三，occupancy 本身是手段，不是目标。它高只说明调度器有足够多的 warp 可以切换，并不自动保证访存效率；反过来，低 occupancy 配合足够的指令级并行，同样能跑满内存带宽。Volkov 的 \textit{Better Performance at Lower Occupancy}（GTC 2010）是这一观点的经典论述。
\end{lookback}

% ================================================================
\section{计时与异步初步}

在前面的练习里 \texttt{cudaDeviceSynchronize} 反复出现以保证计时的正确性。本 Module 主要关注两点——kernel 启动是异步的；以及如何准确计时。

\begin{reading}
Guide 2.5（只读引言）、1.2.3.3。
\end{reading}

\prob{5.1}{EXPERIMENT}[cuda/m5\_async/01\_timing\_trap.cu]
运行实验：

\begin{lstlisting}
cd assignment01/cuda
make run/m5_async/01_timing_trap
\end{lstlisting}

并回答下列问题：
(a)~哪个数值可以当作 kernel 耗时写进报告？(b)~另外两个各具体测的是什么？

\prob{5.2}{CONCEPT}
判断对错，可以顺带补一句理由。
\begin{enumerate}[label=(\alph*)]
  \item 同一个 stream 里的操作按提交顺序执行。
  \item kernel 启动后，host 代码立刻继续往下执行。
  \item unified memory 下，CPU 访问一页正被 GPU 占用的内存，会触发缺页与页迁移。
\end{enumerate}

\begin{lookback}
Module 5 的两道题分量不算重，但在整个认知链条里恰好落在一个关键的位置。CUDA 里的异步并不是某种需要额外开启的特性，它更像是 kernel launch 的默认语义：你把工作提交到某条 stream 上，host 端立刻返回，device 端的执行则在另一条时间线上独立展开。正因如此，“计时”这件事本身变成了一门需要认真对待的事——warmup、同步点插在哪里、选哪种 timer、重复多少次，任何一环处理得马虎，拿到的数字就失去了意义。这次作业里，这些细节已经由写好的测试框架替你包揽了；等到将来你自己搭建 benchmark，就得把这些一一补回来。

再下一步，异步自然会引出重叠：多 stream、数据拷贝与计算的并发、用 CUDA Graph 把整张任务图一次性提交。这些内容会放在后续 session 里展开，这里只需要建立起一个清晰的 mental model：host 只管提交，device 只管执行，两条时间线各自推进。
\end{lookback}

% ================================================================
\section{Tile 视角}

Guide 从 13.x 起把 tile 编程作为与 SIMT 并列的第二种官方模型写进正文。tile 模型描述的是“一个 block 对一块数据做什么”，而 block 内 threads 的分工由编译器决定。此 Module 只做概念铺垫和对照阅读。

\begin{reading}
Guide 1.2.2.3（必读，篇幅不长）、2.4.1--2.4.6（浏览）、2.2（知道 CUDA Python 这条路线存在即可）。
\end{reading}

\prob{6.1}{CONCEPT}
判断对错，可以顺带补一句理由。
\begin{enumerate}[label=(\alph*)]
  \item tile 是显存里的一块可变区域，kernel 通过指针直接改写它。
  \item 对 tile 的一次运算（如两个 tile 相加）由编译器映射到 block 内的多个线程上执行。
  \item tile 模型与 SIMT 模型互斥，一个 CUDA 程序只能选一种。
\end{enumerate}

\prob{6.2}{CONCEPT}
下面是 Guide 2.4.6 的 cuTile Python 向量加法：

\begin{lstlisting}[language=Python]
import cuda.tile as ct

@ct.kernel
def vec_add(a, b, c, TILE: ct.Constant[int]):
    a_view = a.tiled_view((TILE,))
    b_view = b.tiled_view((TILE,))
    c_view = c.tiled_view((TILE,))

    bid = ct.bid(0)
    a_tile = a_view.load((bid,))
    b_tile = b_view.load((bid,))
    c_view.store((bid,), a_tile + b_tile)
\end{lstlisting}

据此补全下表（Triton 一列可以做完 Module 7 再回来填）：

\begin{center}\small
\begin{tabular}{@{}lp{8em}p{9em}p{9em}@{}}
\toprule
 & CUDA SIMT & cuTile & Triton \\
\midrule
并行单位 & block 里的 thread & block & \\
编号 & \texttt{blockIdx} / \texttt{threadIdx} &  & \\
数据分工 & 线程用全局下标来划分数据 &  & \\
边界处理 & \texttt{if} 判断 &  & \\
\bottomrule
\end{tabular}
\end{center}

\prob{6.3}{CONCEPT}
仍看上面这段代码。(a)~“每个线程对应哪个/些元素”由谁决定？(b)~列出一些在 CUDA SIMT 版向量加法里一定会出现、这里完全没体现出的概念。

\begin{lookback}
tile 并非某个 DSL 的发明。NVIDIA 自己把它与 SIMT 并列写进了 Guide 正文，这本身就说明它是一类稳定的编程方式，而不是一次工具选型。

从 SIMT 转向 tile，真正改变的是抽象层的职责划分：线程到数据的映射交给了编译器，tile 的尺寸却仍由用户来定。这条界线不是随意画下的——前者凭形状就可以机械推导，后者必须依赖硬件参数与实测，编译器暂时还无力接手。prob 7.5 的表格会把这一判据推广到四种模型上。

需要明确的是，抽象层的上移并不会抹掉底层的硬件。coalescing、shared memory 的容量、occupancy 这些约束仍然在起作用，只不过不再由你亲手去处理。这也是本作业先用五个模块讲清楚 SIMT，再引入 DSL 的缘故。
\end{lookback}

% ================================================================
\section{TileLang 与 Triton}

此 Module 用 Triton / TileLang 重写一部分之前用 CUDA SIMT 实现的 kernel，重心在 TileLang。
注：Triton 题（7.1、7.2、7.8）不需要 GPU 也能完成正确性部分，运行方式见 \texttt{README.md}；TileLang 题都需要 GPU，在集群上跑。

\begin{reading}
TileLang Language Basics\footnote{\url{https://tilelang.com/programming_guides/language_basics.html}}（建议阅读）；Triton 官方教程 01-vector-add\footnote{\url{https://triton-lang.org/main/getting-started/tutorials/}}；tilelang-puzzles\footnote{\url{https://github.com/tile-ai/tilelang-puzzles}} 。
\end{reading}

\prob{7.1}{FILL-IN}[kernels/vector\_add.py]
请填空：\texttt{kernels/vector\_add.py} ，具体要求见代码注释。测试命令：

\begin{lstlisting}
cd assignment01
uv run pytest tests/test_vector_add.py
\end{lstlisting}

\prob{7.2}{MODIFY}[kernels/fused\_op.py]
按文件开头注释要求修改：\texttt{kernels/fused\_op.py} 。测试命令：

\begin{lstlisting}
cd assignment01
uv run pytest tests/test_fused_op.py
\end{lstlisting}

改完回答——与 Module 2 里改 CUDA kernel 相比，这次的改动主要集中在 kernel 的什么部分？主体代码为什么一行都不用动？

\prob{7.3}{FILL-IN}[kernels/tilelang\_scale\_add.py]
请填空：\texttt{kernels/tilelang\_scale\_add.py} ，具体要求见代码注释。测试命令：

\begin{lstlisting}
cd assignment01
uv sync --extra tilelang
uv run pytest tests/test_tilelang.py -k scale_add
\end{lstlisting}

\prob{7.4}{FILL-IN}[kernels/tilelang\_copy2d.py]
请填空：\texttt{kernels/tilelang\_copy2d.py} ，具体要求见代码注释。测试命令：

\begin{lstlisting}
cd assignment01
uv run pytest tests/test_tilelang.py -k copy2d
\end{lstlisting}

填完想一想：2.6 的四个空——行号、列号、边界保护、grid 尺寸——哪些在这里还有对应？没有对应的那个去哪了？

\prob{7.5}{CONCEPT}
补全下表，每个空填“用户”或“编译器”（二者都涉及的要写清楚各自的范围）。

\begin{center}\small
\begin{tabular}{@{}lcccc@{}}
\toprule
谁负责 & CUDA SIMT & cuTile & Triton & TileLang \\
\midrule
线程到数据的映射 & 用户 & & & \\
边界处理 & 用户 & & & \\
tile / block 尺寸的选择 & 用户 & & & \\
block 内同步 & 用户 & & & \\
\bottomrule
\end{tabular}
\end{center}

\prob{7.6}{FILL-IN}[kernels/tilelang\_matmul.py]
请填空：\texttt{kernels/tilelang\_matmul.py} ，具体要求见代码注释。测试命令：

\begin{lstlisting}
cd assignment01
uv run pytest tests/test_tilelang.py -k matmul
\end{lstlisting}

填完回答——这五个空涉及到了 shared memory、寄存器 tile、流水，而 Triton 版 matmul（Bonus 的 \texttt{kernels/matmul\_triton.py}）并未被显式指定，为什么？

\begin{capstone}{prob 7.7（FROM-SCRATCH）：softmax in TileLang \hfill {\footnotesize\ttfamily kernels/tilelang\_softmax.py}}
在 \texttt{kernels/tilelang\_softmax.py} 里从零实现行 softmax，具体要求见文件开头的 docstring。
\end{capstone}

\begin{lstlisting}
cd assignment01
uv run pytest tests/test_tilelang_softmax.py
\end{lstlisting}

\prob[Optional]{7.8}{FROM-SCRATCH}[kernels/softmax.py]
用 Triton 重写 prob 7.7，此题可以不用 GPU。写完后可以试试对比一下此题与 7.7 ：归约、边界处理、按形状编译，二者各自是由谁处理的（用户显式处理或者编译器隐式处理）？

测试命令：

\begin{lstlisting}
cd assignment01
uv run pytest tests/test_softmax.py
\end{lstlisting}

\begin{lookback}
本模块把前面写过的 kernel 用 TileLang 和 Triton 各重写了一遍。两种写法的区别，基本都收在 prob 7.5 那张对比表里：线程到数据的映射归谁管，边界归谁处理，tile 尺寸归谁选定，同步又归谁插入。DSL 的价值并不在于代码更短，而在于它把前两项交给编译器，同时把后两项——那些必须靠实测才能定下来的旋钮——继续留给你。

抽象当然有代价。Bonus 部分的一组数字提醒你，表达能力和实现成熟度需要分开评估：同一段 TileLang 换个版本，跑出的结果就可能不一样。更关键的是，一旦编译器接手的那部分出了问题——layout 不合法、精度对不上、边界被悄悄吃掉——你的排查路径最终还是得退回 SIMT 这一层，去看它究竟生成了什么。DSL 是构筑在 CUDA 之上的一层，从来不是它的替代品。
\end{lookback}

% ================================================================
\section{平台与编译}

\begin{reading}
Guide 1.3 全章（必读）、2.1.1、2.7（浏览）。
\end{reading}

\prob{8.1}{CONCEPT}
判断对错，可以顺带补一句理由。
\begin{enumerate}[label=(\alph*)]
  \item PTX 是 GPU 直接执行的机器码。
  \item 只嵌入了 \texttt{sm\_70} SASS 的可执行文件，能在 compute capability 9.0 的卡上运行。
  \item 一个 fatbin 可以同时携带多个架构的 SASS 和 PTX。
  \item JIT 编译由驱动在运行时完成。
\end{enumerate}

\prob[Optional]{8.2}{EXPERIMENT}[cuda/m0\_env/01\_hello.cu]
两个编译实验，对象是 Module 0 的 \texttt{01\_hello.cu}。

(a)~生成只含 \texttt{sm\_90} SASS 的可执行文件并运行（如果你的卡本身就是 compute capability 9.0，先把 \texttt{ARCH\_HIGH} 调成 100 或更高再 make），记录报错信息。

\begin{lstlisting}
cd assignment01/cuda
make sassonly/m0_env/01_hello
./bin/m0_env/01_hello_sassonly
\end{lstlisting}

(b)~生成只含 \texttt{compute\_75} PTX 的版本（CUDA 13 起 \texttt{compute\_70} 已被移除，Makefile 默认取 75）并运行。能正常运行吗？PTX 是在什么时候、由谁编译成这块卡的机器码的？

\begin{lstlisting}
cd assignment01/cuda
make ptxonly/m0_env/01_hello
./bin/m0_env/01_hello_ptxonly
\end{lstlisting}

\prob[Optional]{8.3}{CONCEPT}
请简单说明 Runtime API 与 Driver API 各自的定位。\texttt{cudaMalloc} 属于哪个？

\begin{lookback}
Module 0 里那个“能跑就行”的 hello，现在被拆开重新看过：源码经 nvcc 分成 host 和 device 两路，device 端产出 PTX 与 SASS，一并打包进 fatbin；运行时 driver 再从中挑选合适的版本，或者对 PTX 做 JIT。

这一模块的实用意义在交付环节。“编译通过”和“能在目标卡上跑起来”是两回事，后者是部署中最常撞上的墙。当你把一个 GPU 程序交付给别人，至少需要说清楚：它支持哪些 compute capability，fatbin 里包含了哪些架构的 SASS，有没有 PTX 兜底，以及首次运行的 JIT 开销落在哪里。写 kernel 与发布软件是两种视角，而这里正是二者的交界。
\end{lookback}

% ================================================================
\clearpage
\section*{Bonus（EXPERIMENT · Optional）：matmul}
\addcontentsline{toc}{section}{Bonus（EXPERIMENT · Optional）：matmul}

调参并比较不同 matmul 实现的性能差距。

\begin{enumerate}[label=(\alph*)]
  \item Naive CUDA：\texttt{cuda/bonus/matmul.cu}，用 \texttt{-DBS} 取 8、16、32 各跑一遍，记录实测数据。
\begin{lstlisting}
cd assignment01/cuda
make run/bonus/matmul
nvcc -O2 -std=c++17 -I. -arch=native -DBS=8 -o bin/bonus/matmul_bs8 bonus/matmul.cu && ./bin/bonus/matmul_bs8
\end{lstlisting}
  \item Tiled Triton：\texttt{kernels/matmul\_triton.py}（多组 tile 配置 + cuBLAS 对照），记录实测数据。
\begin{lstlisting}
cd assignment01
uv run python -c "from kernels.matmul_triton import bench; bench()"
\end{lstlisting}
  \item TileLang：\texttt{kernels/tilelang\_matmul.py}（即 prob 7.6 补全的成品），调 \texttt{bench()} 的 \texttt{block\_M / block\_N / block\_K / num\_stages}，记录实测数据。
\begin{lstlisting}
cd assignment01
uv sync --extra tilelang
uv run python -c "from kernels.tilelang_matmul import bench; bench()"
\end{lstlisting}
\end{enumerate}

试分析性能差距的原因，特别是：TileLang 的控制粒度比 Triton 更细，为什么在这组测试里反而略慢？

注：A100 实测参考——naive CUDA 约 2.4 TFLOPS，Triton 约 125 TFLOPS，TileLang 约 118 TFLOPS，cuBLAS 约 173 TFLOPS，而 A100 fp16 tensor core 的标称上限是 312 TFLOPS。Triton 与 TileLang 的数字只是各自 \texttt{bench()} 里那几组配置中的最优，不代表工具的性能上限；tilelang 本身也还在快速迭代（本仓库固定在 0.1.12），数字会随版本变化。
\begin{lookback}
回顾这份作业，主线可以看作三次视角的切换：单线程到 warp（Module 1--3）、计算到数据搬运（Module 4--5）、线程到 tile（Module 6--7）。走完这一遍，你写下的 kernel 应当能保证正确，且不至于慢得离谱。但距离“快”字仍然有一段差距，Bonus 里那组数字恰好把它量化了出来——naive 实现和 cuBLAS 之间，隔着 tensor core、asynchronous copy、software pipelining 以及 profiler-driven tuning，这些内容会在后续的 session 中展开。
\end{lookback}
\end{document}
